%! Least Squares Approximations — Unit 4, Lesson 4.3 — Homework
\documentclass[10pt]{article}
\usepackage{linalg-article}
\usepackage{linalg-boxes}
\newcommand{\TallMath}[1]{$\displaystyle #1\rule[-1.4em]{0pt}{3.2em}$}
\newcommand{\xx}{\mathbf{x}}
\newcommand{\bb}{\mathbf{b}}
\newcommand{\pp}{\mathbf{p}}
\newcommand{\ee}{\mathbf{e}}
\newcommand{\zero}{\mathbf{0}}
\newcommand{\T}{\mathsf{T}}

\begin{document}

\pageheader{Unit 4, Lesson 4.3}{Homework}
\namedateperiod

\vspace{0.12in}

\begin{notesbox}{Practice}
To fit $y=C+Dt$ to data: build $A$ (a column of $1$s and a column of the $t$-values) and $\bb$ (the
$y$-values), then solve the normal equations $A^{\T}A\hat{\xx}=A^{\T}\bb$. The fitted values are
$\pp=A\hat{\xx}$ and the residuals are $\ee=\bb-\pp$. Show your dot products.

\begin{enumerate}[leftmargin=1.9em, itemsep=11pt]
  \item \textbf{Best flat line.} Fit a horizontal line $y=C$ to $\bb=\begin{bmatrix} 2\\5\\8 \end{bmatrix}$
        (so $A=\begin{bmatrix} 1\\1\\1 \end{bmatrix}$). Find $\hat{C}$, the fitted values $\pp$, and the
        residuals $\ee$. What familiar quantity is $\hat{C}$? \writelines{2}
  \item \textbf{Line through three points.} Fit $y=C+Dt$ to $(0,1),(1,0),(2,5)$. Find $A^{\T}A$,
        $A^{\T}\bb$, then $\hat{\xx}$, the line, $\pp$, and $\ee$; check $\ee$ is perpendicular to both
        columns. \writelines{3}
  \item \textbf{Line through four points, then predict.} A plant's height (cm) at weeks $t=0,1,2,3$ is
        $\bb=\begin{bmatrix} 4\\4\\6\\10 \end{bmatrix}$.
    \begin{enumerate}[label=(\alph*), leftmargin=1.8em, itemsep=5pt]
      \item Solve the normal equations for $\hat{\xx}=(C,D)$ and state the best-fit line. \vspace{1.2cm}
      \item Give $\pp$ and $\ee$, then predict the height at week $t=4$. \vspace{1.0cm}
    \end{enumerate}
  \item \textbf{Explain.} A classmate says ``the best line should pass through all the points.'' In two or
        three sentences, explain why that is usually impossible and what least squares does instead.
        \writelines{2}
\end{enumerate}
\end{notesbox}

\begin{extensionbox}
\textbf{Where the residual lives.} Use your answers from Problem~3
($A=\begin{bsmallmatrix} 1&0\\1&1\\1&2\\1&3 \end{bsmallmatrix}$ and $\ee$).
\begin{enumerate}[label=(\alph*), leftmargin=1.8em, itemsep=5pt]
  \item Check $A^{\T}\ee=\zero$. Which of the four subspaces holds $\ee$, and which space is it the
        orthogonal complement of (Lesson~4.1)?
  \item Because the first column of $A$ is all $1$s, one entry of $A^{\T}\ee=\zero$ says $\sum e_i=0$. Show
        this for your $\ee$, and explain why the residuals of a best-fit line always add to zero.
\end{enumerate}
\writelines{3}
\end{extensionbox}

\begin{spiralbox}
\textbf{Coming next --- Lesson 4.4: Orthogonal Matrices \& Gram--Schmidt.} Solving $A^{\T}A\hat{\xx}=A^{\T}\bb$
is easy when the columns of $A$ are \emph{perpendicular unit vectors} --- then $A^{\T}A=I$ and the
projection is immediate. Next lesson builds such orthonormal columns and the matrix $Q$ that carries them.
\end{spiralbox}

\end{document}
